\documentclass[12pt,a4paper]{article}
\usepackage[margin=15mm]{geometry}
\usepackage{amsmath,amsthm,amssymb,amsfonts}
\usepackage{enumitem}
\usepackage{xcolor}
\usepackage{hyperref}
\hypersetup{colorlinks=true,linkcolor=blue!70!black,urlcolor=blue!80!black}

\theoremstyle{definition}
\newtheorem{exercise}{Problem}
\newtheorem*{solution*}{Solution}
\usepackage{parskip}
\usepackage{etoolbox}
\AtBeginEnvironment{exercise}{\vspace{\parskip}}
\AtBeginEnvironment{solution*}{\vspace{\parskip}}

\newcommand\expc{\mathsf{E}}
\newcommand\prb{\mathsf{P}}
\DeclareMathOperator\var{var}
\DeclareMathOperator\cov{cov}
\DeclareMathOperator\tr{tr}
\DeclareMathOperator\rank{rank}
\DeclareMathOperator\diag{diag}
\DeclareMathOperator{\avec}{vec}
\DeclareMathOperator{\avech}{vech}
\newcommand{\dd}{\mathrm{d}}
\newcommand{\SST}{\mathrm{SST}}
\newcommand{\SSR}{\mathrm{SSR}}
\newcommand{\SSE}{\mathrm{SSE}}
\newcommand{\MSR}{\mathrm{MSR}}
\newcommand{\MSE}{\mathrm{MSE}}
\newcommand{\SE}{\mathrm{SE}}
\newcommand{\bX}{\mathbf{X}}
\newcommand{\bY}{\mathbf{Y}}
\newcommand{\bx}{\mathbf{x}}
\def\bb{\boldsymbol{\beta}}
\newcommand{\be}{\boldsymbol{\varepsilon}}
\newcommand{\bI}{\mathbf{I}}
\newcommand{\bH}{\mathbf{H}}
\newcommand{\bM}{\mathbf{M}}
\newcommand{\bR}{\mathbf{R}}
\newcommand{\br}{\mathbf{r}}
\def\bc{\mathbf{c}}
\newcommand{\bone}{\mathbf{1}}
\newcommand{\bzero}{\mathbf{0}}
\newcommand{\bz}{\mathbf{z}}
\newcommand{\bA}{\mathbf{A}}
\newcommand{\bC}{\mathbf{C}}
\newcommand{\bD}{\mathbf{D}}
\newcommand{\bP}{\mathbf{P}}
\newcommand{\ba}{\mathbf{a}}
\newcommand{\bv}{\mathbf{v}}
\newcommand{\bw}{\mathbf{w}}
\newcommand{\bs}{\mathbf{s}}
\newcommand{\bt}{\mathbf{t}}
\newcommand{\Real}{\mathbb{R}}
\newcommand{\calC}{\mathcal{C}}

\title{\textbf{Multivariate Analysis: Midterm Test Bank}\\[2mm]\large Matrix Calculus \& Multiple Linear Models}
\author{}
\date{}

\begin{document}
\maketitle

\medskip\noindent\textbf{Conventions.}\; $\dd X$ is the differential of $X$; $\dd^2 X = \dd(\dd X)$. For scalar $f$: $Df(x) = \partial f/\partial x'$ (row vector); $Hf(x) = \partial^2 f/\partial x\,\partial x'$ (Hessian). For matrix $F(X)$: $DF(X) = \partial\avec F/\partial(\avec X)'$. Let $e_i$ denote the $i$-th standard basis vector.

\medskip\noindent\textbf{Identities you may cite without proof.}
\begin{enumerate}[label=(I\arabic*),nosep]
\item $\tr A' = \tr A$; \; $\dd(\tr X) = \tr\dd X$; \; $\dd A = 0$ for constant $A$.
\item $\dd(XY) = (\dd X)Y + X\,\dd Y$; \; transpose of a scalar: $(\dd x)'Ax = x'A'\dd x$.
\item $\dd^2 x = 0$ when $x$ is the independent variable.
\end{enumerate}

%=============================================================================
\bigskip
\begin{center}
\rule{0.6\textwidth}{0.4pt}\\[2mm]
{\Large\bfseries Part I: Matrix Calculus}\\[1mm]
\rule{0.6\textwidth}{0.4pt}
\end{center}

\subsection*{Tier A: Reproduce the Derivation}

%--- MC-1 ---
\begin{exercise}[Trace cycling]\label{mc:A1}
Prove that for any two matrices $A$ and $B$ of the same order, $\tr A'B = \tr BA'$.
\end{exercise}

\begin{solution*}
By definition of the trace and matrix product:
\[
\tr A'B = \sum_j (A'B)_{jj} = \sum_j \sum_i (A')_{ji}\, b_{ij} = \sum_j \sum_i a_{ij}\, b_{ij}.
\]
Interchanging the order of summation:
\[
= \sum_i \sum_j b_{ij}\, a_{ij} = \sum_i (BA')_{ii} = \tr BA'. \tag*{\qed}
\]
\end{solution*}

%--- MC-2 ---
\begin{exercise}[Quadratic form equals zero for all $x$]\label{mc:A2}
Prove: (a)~$x'Ax = 0$ for every $x$ if and only if $A$ is skew-symmetric ($A' = -A$). (b)~If additionally $A = A'$, then $A = 0$.
\end{exercise}

\begin{solution*}
\textbf{(a) ($\Leftarrow$):} If $A' = -A$, then $x'Ax$ is a scalar, so $x'Ax = (x'Ax)' = x'A'x = -x'Ax$. Hence $2x'Ax = 0$, i.e., $x'Ax = 0$ for all $x$.

\textbf{($\Rightarrow$):} Assume $x'Ax = 0$ for every $x$. Setting $x = e_i$ (the $i$-th standard basis vector): $e_i'Ae_i = a_{ii} = 0$ for all $i$. Setting $x = e_i + e_j$ ($i \neq j$):
\[
(e_i + e_j)'A(e_i + e_j) = a_{ii} + a_{ij} + a_{ji} + a_{jj} = 0.
\]
Since $a_{ii} = a_{jj} = 0$, we get $a_{ij} + a_{ji} = 0$ for all $i \neq j$. Combined with $a_{ii} = 0$, this gives $A' = -A$.

\textbf{(b)} From part (a), $x'Ax = 0$ for all $x$ implies $A = -A'$. Combined with $A = A'$: $A = -A$, so $2A = 0$ and $A = 0$. \qed
\end{solution*}

%--- MC-4 ---
\begin{exercise}[$\avec(ABC) = (C' \otimes A)\avec B$]\label{mc:A4}
Prove that for any compatible matrices $A$, $B$, $C$: $\avec(ABC) = (C' \otimes A)\avec B$.
\end{exercise}

\begin{solution*}
Let $b_j$ denote the $j$-th column of $B$ and $e_j$ the $j$-th standard basis vector (of appropriate dimension). Then $B = \sum_{j} b_j e_j'$, so
\[
ABC = \sum_{j} A(b_j e_j')C = \sum_{j} (Ab_j)(C'e_j)'.
\]
Applying $\avec(\cdot)$ and using $\avec(ab') = b \otimes a$:
\begin{align*}
\avec(ABC) &= \sum_{j} \avec\bigl((Ab_j)(C'e_j)'\bigr) = \sum_{j} (C'e_j) \otimes (Ab_j) \\
&= \sum_{j} (C' \otimes A)(e_j \otimes b_j) = (C' \otimes A) \sum_{j} \avec(b_j e_j') = (C' \otimes A)\avec B. \tag*{\qed}
\end{align*}
\end{solution*}

%--- MC-5 ---
\begin{exercise}[Differential of the determinant]\label{mc:A5}
Prove that for nonsingular $X$: $\dd|X| = |X|\,\tr X^{-1}\dd X$, and hence $\dd\log|X| = \tr X^{-1}\dd X$ when $|X| > 0$.
\end{exercise}

\begin{solution*}
By the chain rule for scalar functions:
\[
\dd|X| = \sum_{i,j} \frac{\partial|X|}{\partial x_{ij}}\,\dd x_{ij} = \sum_{i,j} C_{ij}\,\dd x_{ij}.
\]
Substituting $C_{ij} = |X|(X^{-1})_{ji}$:
\begin{align*}
\dd|X| &= |X| \sum_{i,j} (X^{-1})_{ji}\,\dd x_{ij} = |X| \sum_{j} \sum_{i} (X^{-1})_{ji}\,\dd x_{ij} = |X| \sum_{j} (X^{-1}\dd X)_{jj} = |X|\,\tr X^{-1}\dd X.
\end{align*}
For $|X| > 0$: $\dd\log|X| = \dd|X|/|X| = \tr X^{-1}\dd X$. \qed
\end{solution*}

%--- MC-6 ---
\begin{exercise}[Differential of the inverse]\label{mc:A6}
Prove that for nonsingular $X$: $\dd X^{-1} = -X^{-1}(\dd X)X^{-1}$.
\end{exercise}

\begin{solution*}
Start from the identity $X^{-1}X = I$. Differentiating both sides using the product rule:
\[
(\dd X^{-1})X + X^{-1}(\dd X) = \dd I = 0.
\]
Solving for $\dd X^{-1}$:
\[
(\dd X^{-1})X = -X^{-1}(\dd X).
\]
Postmultiplying both sides by $X^{-1}$:
\[
\dd X^{-1} = -X^{-1}(\dd X) X^{-1}. \tag*{\qed}
\]
\end{solution*}

\subsection*{Tier B: Apply the Theory}

%--- MC-8 ---
\begin{exercise}[$\avec(ab') = b \otimes a$]\label{mc:B1}
Let $a$ be $m \times 1$ and $b$ be $n \times 1$. Show that $\avec(ab') = b \otimes a$.
\end{exercise}

\begin{solution*}
The matrix $ab'$ is $m \times n$. Its $j$-th column is $b_j a$ (the scalar $b_j$ times the vector $a$). The $\avec$ operator stacks the columns:
\[
\avec(ab') = \begin{pmatrix} b_1 a \\ b_2 a \\ \vdots \\ b_n a \end{pmatrix} = b \otimes a,
\]
by the definition of the Kronecker product (each entry $b_j$ of $b$ multiplies the vector $a$). \qed
\end{solution*}

%--- MC-8 ---
\begin{exercise}[Derivatives of $a'x$ and $x'Ax$ via differentials]\label{mc:B2}
Let $a$ be an $n \times 1$ constant vector and $A$ an $n \times n$ constant matrix.\\
(a) Using differentials, show that $\partial(a'x)/\partial x' = a'$.\\
(b) Using differentials, show that $\partial(x'Ax)/\partial x' = x'(A + A')$.\\
(c) When $A$ is symmetric, simplify (b).
\end{exercise}

\begin{solution*}
\textbf{(a)} Since $a$ is constant, $\dd a = 0$:
\[
\dd(a'x) = (\dd a)'x + a'(\dd x) = 0 + a'\dd x = a'\dd x.
\]
By the identification theorem ($\dd f = c'\dd x \Rightarrow \partial f/\partial x' = c'$), the derivative is $a'$.

\textbf{(b)} Since $A$ is constant, $\dd A = 0$:
\[
\dd(x'Ax) = (\dd x)'Ax + x'(\dd A)x + x'A(\dd x) = (\dd x)'Ax + x'A\,\dd x.
\]
The first term is a $1 \times 1$ scalar, so it equals its own transpose: $(\dd x)'Ax = ((\dd x)'Ax)' = x'A'\dd x$. Hence:
\[
\dd(x'Ax) = x'A'\dd x + x'A\,\dd x = x'(A + A')\dd x.
\]
By the identification theorem, $\partial(x'Ax)/\partial x' = x'(A+A')$.

\textbf{(c)} When $A = A'$: $A + A' = 2A$, so the derivative is $2x'A$. \qed
\end{solution*}

%--- MC-9 ---
\begin{exercise}[Derivative of $\tr X'AX$]\label{mc:B3}
Let $A$ be a constant $m \times m$ matrix and $X$ an $m \times n$ matrix variable. Show that $Df(X) = (\avec((A+A')X))'$ where $f(X) = \tr X'AX$.
\end{exercise}

\begin{solution*}
Using $\dd(\tr M) = \tr(\dd M)$ and the product rule:
\begin{align*}
\dd f(X) &= \dd(\tr X'AX) = \tr\dd(X'AX) = \tr\bigl[(\dd X)'AX + X'A(\dd X)\bigr].
\end{align*}
For the first term, use $\tr M = \tr M'$:
\[
\tr(\dd X)'AX = \tr\bigl((\dd X)'AX\bigr)' = \tr(X'A'\dd X).
\]
Combining:
\[
\dd f = \tr X'A'\dd X + \tr X'A\,\dd X = \tr X'(A+A')\dd X = \tr C'\dd X,
\]
where $C = (A+A')X$. By the identification theorem for matrices ($\dd f = \tr C'\dd X \Leftrightarrow Df = (\avec C)'$):
\[
Df(X) = \bigl(\avec\bigl((A+A')X\bigr)\bigr)'. \tag*{\qed}
\]
\end{solution*}

%--- MC-10 ---
\begin{exercise}[Derivative of $\log|X'X|$]\label{mc:B4}
Let $X$ be an $m \times n$ matrix with full column rank. Show that $Df(X) = 2(\avec(X(X'X)^{-1}))'$ where $f(X) = \log|X'X|$.
\end{exercise}

\begin{solution*}
Let $Y = X'X$. Using $\dd\log|Y| = \tr Y^{-1}\dd Y$:
\[
\dd f = \tr(X'X)^{-1}\dd(X'X) = \tr(X'X)^{-1}\bigl[(\dd X)'X + X'\dd X\bigr].
\]
For the first term: $\tr(X'X)^{-1}(\dd X)'X$ is a scalar, so it equals its transpose. Computing:
\[
\bigl[(X'X)^{-1}(\dd X)'X\bigr]' = X'(\dd X)(X'X)^{-1},
\]
so $\tr(X'X)^{-1}(\dd X)'X = \tr X'(\dd X)(X'X)^{-1} = \tr(X'X)^{-1}X'\dd X$ (by trace cycling). Hence:
\[
\dd f = 2\tr(X'X)^{-1}X'\dd X = 2\tr C'\dd X,
\]
where $C = X(X'X)^{-1}$. By the identification theorem: $Df(X) = 2(\avec C)'$. \qed
\end{solution*}

%--- MC-11 ---
\begin{exercise}[Derivative of $\tr X^k$]\label{mc:B5}
Let $f_k(X) = \tr X^k$ ($k = 1, 2, \ldots$), $X$ square. Find $Df_k(X)$.
\end{exercise}

\begin{solution*}
Using the product rule for $\dd(X^k) = \dd(X \cdot X^{k-1})$:
\[
\dd(X^k) = (\dd X)X^{k-1} + X(\dd X)X^{k-2} + \cdots + X^{k-1}(\dd X).
\]
Taking the trace and applying trace cycling ($\tr X^j(\dd X)X^{k-1-j} = \tr X^{k-1}\dd X$) to each of the $k$ terms:
\[
\dd f_k = k\,\tr X^{k-1}\dd X.
\]
Writing $\tr X^{k-1}\dd X = \tr((X')^{k-1})'\dd X = (\avec (X')^{k-1})'\dd\avec X$ (by $\tr A'B = (\avec A)'\avec B$), we obtain $Df_k(X) = k(\avec (X')^{k-1})'$.

Special cases: $D(\tr X) = (\avec I)'$; \; $D(\tr X^2) = 2(\avec X')'$. \qed
\end{solution*}

%--- MC-12 ---
\begin{exercise}[Derivative of $F(X) = AX^{-1}B$]\label{mc:B6}
$X$ nonsingular, $A$ and $B$ constant. Show $DF(X) = -(X^{-1}B)' \otimes (AX^{-1})$.
\end{exercise}

\begin{solution*}
Using $\dd X^{-1} = -X^{-1}(\dd X)X^{-1}$ (Problem~\ref{mc:A6}):
\[
\dd F(X) = A(\dd X^{-1})B = -AX^{-1}(\dd X)X^{-1}B.
\]
This has the form $M_1(\dd X)M_3$ with $M_1 = -AX^{-1}$ and $M_3 = X^{-1}B$. Vectorizing using $\avec(M_1(\dd X)M_3) = (M_3' \otimes M_1)\dd\avec X$:
\[
\dd\avec F = -\bigl((X^{-1}B)' \otimes (AX^{-1})\bigr)\dd\avec X.
\]
By the identification theorem for matrices: $DF(X) = -(X^{-1}B)' \otimes (AX^{-1})$. \qed
\end{solution*}


%=============================================================================
\bigskip
\begin{center}
\rule{0.6\textwidth}{0.4pt}\\[2mm]
{\Large\bfseries Part II: Multiple Linear Regression}\\[1mm]
\rule{0.6\textwidth}{0.4pt}
\end{center}

\medskip
\noindent\textbf{Model.}\; $\bY = \bX\bb + \be$, $\bY$ is $n \times 1$, $\bX$ is $n \times (k+1)$, $\rank(\bX) = k+1$, $n > k+1$.

\noindent\textbf{Notation.}\; $\bH = \bX(\bX'\bX)^{-1}\bX'$, $\bM = \bI - \bH$, $\hat{\bb} = (\bX'\bX)^{-1}\bX'\bY$, $\mathbf{e} = \bM\bY$, $s^2 = \MSE = \SSE/(n{-}k{-}1)$.

\noindent\textbf{Assumptions.}\; (A1)~$\expc[\bY|\bX]=\bX\bb$; (A2)~$\expc[\be|\bX]=\bzero$; (A3)~$\var(\be|\bX)=\sigma^2\bI$; (A4)~$\rank(\bX)=k{+}1$, $n>k{+}1$; (A5)~$\be|\bX \sim N(\bzero,\sigma^2\bI)$.

\subsection*{Tier A: Reproduce the Derivation}

%--- LM-A1 (clm Lemma 2.2): trace = rank for sym idempotent ---
\begin{exercise}[Trace equals rank for symmetric idempotent matrices]\label{prob:A1}
Let $\bA \in \Real^{n\times n}$ be symmetric and idempotent ($\bA^2 = \bA = \bA'$).
\begin{enumerate}[label=(\alph*),nosep]
\item Show that every eigenvalue of $\bA$ is $0$ or $1$.
\item Conclude $\tr(\bA) = \rank(\bA)$.
\item Verify the formula on $\bH$ and $\bM$: $\tr(\bH) = k+1$ and $\tr(\bM) = n-k-1$.
\end{enumerate}
This trace$=$rank identity is invoked repeatedly throughout the regression theory (degrees of freedom for $\SSE$, $\SSR$, $F$ tests, Cochran's theorem).
\end{exercise}

\begin{solution*}
\textbf{(a)} Let $\bA\bv = \lambda\bv$ with $\bv \neq \bzero$. Idempotency gives $\bA^2\bv = \lambda^2\bv$, but also $\bA^2\bv = \bA\bv = \lambda\bv$. Hence $\lambda^2 = \lambda$, so $\lambda \in \{0,1\}$.

\textbf{(b)} Symmetry implies $\bA$ is orthogonally diagonalizable: $\bA = \bP\boldsymbol\Lambda\bP'$ with $\bP$ orthogonal and $\boldsymbol\Lambda = \diag(\lambda_1,\ldots,\lambda_n)$. By (a), each $\lambda_i \in \{0,1\}$. Then
\[
\tr(\bA) = \tr(\bP\boldsymbol\Lambda\bP') = \tr(\boldsymbol\Lambda\bP'\bP) = \tr(\boldsymbol\Lambda) = \#\{i : \lambda_i = 1\} = \rank(\boldsymbol\Lambda) = \rank(\bA).
\]

\textbf{(c)} For $\bH = \bX(\bX'\bX)^{-1}\bX'$:
\[
\tr(\bH) = \tr\bigl((\bX'\bX)^{-1}\bX'\bX\bigr) = \tr(\bI_{k+1}) = k+1.
\]
Then $\tr(\bM) = \tr(\bI_n) - \tr(\bH) = n-k-1$. \qed
\end{solution*}

%--- LM-A2 (clm Theorem 2.1): OLS Derivation ---
\begin{exercise}[OLS Derivation]\label{prob:A2}
Derive $\hat{\bb} = (\bX'\bX)^{-1}\bX'\bY$ by minimizing $\|\bY - \bX\bb\|^2$. Verify the second-order condition.
\end{exercise}

\begin{solution*}
\textbf{Step 1: Expand $Q(\bb)$.}
\[
Q(\bb) = (\bY - \bX\bb)'(\bY - \bX\bb) = \bY'\bY - 2\bb'\bX'\bY + \bb'\bX'\bX\bb,
\]
using $\bY'\bX\bb = \bb'\bX'\bY$ (a scalar equals its transpose).

\textbf{Step 2: First-order condition.} Differentiating w.r.t.\ $\bb'$ (with $\bA = \bX'\bX$ symmetric):
\[
\frac{\partial Q}{\partial \bb'} = -2\bY'\bX + 2\bb'\bX'\bX = \bzero'.
\]
Transposing yields the normal equation $\bX'\bX\hat{\bb} = \bX'\bY$.

\textbf{Step 3: Solve.} Under (A4), $\rank(\bX) = k+1$, hence $\rank(\bX'\bX) = \rank(\bX) = k+1$ (since $\bX\bv = \bzero \Leftrightarrow \bX'\bX\bv = \bzero$, the two matrices share a null space). So $\bX'\bX$ is invertible:
\[
\hat{\bb} = (\bX'\bX)^{-1}\bX'\bY.
\]

\textbf{Step 4: Second-order condition.} The Hessian is $\partial^2 Q/(\partial\bb\,\partial\bb') = 2\bX'\bX$. For any $\bv \neq \bzero$, $\bv'\bX'\bX\bv = \|\bX\bv\|^2 > 0$ by (A4). Hence $\bX'\bX \succ \bzero$, so $\hat{\bb}$ is the global minimum. \qed
\end{solution*}

%--- LM-A3 (clm Theorem 2.3): Properties of H, M ---
\begin{exercise}[Properties of $\bH$ and $\bM$]\label{prob:A3}
Let $\bH = \bX(\bX'\bX)^{-1}\bX'$ and $\bM = \bI - \bH$.
\begin{enumerate}[label=(\alph*),nosep]
\item Show $\bH$ and $\bM$ are symmetric and idempotent.
\item Show $\bH\bX = \bX$, $\bM\bX = \bzero$, and $\bH\bM = \bM\bH = \bzero$.
\item Show $\bH$ projects $\Real^n$ onto $\calC(\bX)$ and $\bM$ projects onto $\calC(\bX)^\perp$. (Recall $\calC(\bX)$ is the column space of $\bX$.)
\item Use Problem~\ref{prob:A1} to find $\rank(\bH)$ and $\rank(\bM)$.
\end{enumerate}
\end{exercise}

\begin{solution*}
\textbf{(a)} Symmetry: $\bH' = (\bX(\bX'\bX)^{-1}\bX')' = \bX((\bX'\bX)^{-1})'\bX' = \bX(\bX'\bX)^{-1}\bX' = \bH$ (since $(\bX'\bX)^{-1}$ is symmetric as the inverse of a symmetric matrix). Idempotent:
\[
\bH^2 = \bX(\bX'\bX)^{-1}\bX'\bX(\bX'\bX)^{-1}\bX' = \bX(\bX'\bX)^{-1}\bX' = \bH.
\]
Then $\bM' = \bI' - \bH' = \bM$ and $\bM^2 = (\bI-\bH)^2 = \bI - 2\bH + \bH^2 = \bI - \bH = \bM$.

\textbf{(b)} $\bH\bX = \bX(\bX'\bX)^{-1}\bX'\bX = \bX$, hence $\bM\bX = \bX - \bH\bX = \bzero$. Then $\bH\bM = \bH(\bI-\bH) = \bH - \bH^2 = \bzero$, and $\bM\bH = (\bI-\bH)\bH = \bzero$ by symmetry of the previous identity (or direct computation).

\textbf{(c)} For $\bv \in \Real^n$, $\bH\bv = \bX(\bX'\bX)^{-1}\bX'\bv \in \calC(\bX)$, and $\bH\bv = \bv$ iff $\bv \in \calC(\bX)$ (forward: by (b) any $\bv = \bX\bc$ satisfies $\bH\bv = \bX\bc = \bv$; reverse: if $\bH\bv = \bv$ then $\bv \in \mathrm{Im}(\bH) \subseteq \calC(\bX)$). So $\bH$ is the orthogonal projector onto $\calC(\bX)$. Then $\bM = \bI - \bH$ projects onto the orthogonal complement.

\textbf{(d)} By Problem~\ref{prob:A1}, $\rank(\bH) = \tr(\bH) = k+1$ and $\rank(\bM) = \tr(\bM) = n-k-1$. \qed
\end{solution*}

%--- LM-A4 (clm Theorem 2.4): Algebraic properties of OLS residuals ---
\begin{exercise}[Algebraic properties of OLS residuals]\label{prob:A4}
Assume the model includes an intercept (the first column of $\bX$ is $\bone$). Let $\mathbf{e} = \bY - \bX\hat{\bb}$ and $\hat{\bY} = \bX\hat{\bb}$.
\begin{enumerate}[label=(\alph*),nosep]
\item Derive the normal equations $\bX'\mathbf{e} = \bzero$.
\item Show $\bone'\mathbf{e} = 0$, hence $\bar{e} = 0$.
\item Show $\hat{\bY}'\mathbf{e} = 0$ and $\bY'\mathbf{e} = \mathbf{e}'\mathbf{e}$.
\item Show $\overline{\hat{Y}} = \bar Y$.
\end{enumerate}
\end{exercise}

\begin{solution*}
\textbf{(a)} The first-order condition for OLS is $\partial\|\bY-\bX\bb\|^2/\partial\bb = -2\bX'(\bY - \bX\hat{\bb}) = \bzero$, equivalently $\bX'\mathbf{e} = \bzero$.

\textbf{(b)} The first row of $\bX'\mathbf{e} = \bzero$ reads $\bone'\mathbf{e} = \sum_i e_i = 0$, hence $\bar e = 0$.

\textbf{(c)} $\hat{\bY}'\mathbf{e} = (\bX\hat{\bb})'\mathbf{e} = \hat{\bb}'\bX'\mathbf{e} = \bzero$. Then $\bY'\mathbf{e} = (\hat{\bY} + \mathbf{e})'\mathbf{e} = 0 + \mathbf{e}'\mathbf{e}$.

\textbf{(d)} From $\bY = \hat{\bY} + \mathbf{e}$ summing components and dividing by $n$: $\bar Y = \overline{\hat{Y}} + \bar e = \overline{\hat{Y}}$ by (b). \qed
\end{solution*}

%--- LM-A5 (clm Lemma 3.1): Variance under linear transformation ---
\begin{exercise}[Variance under linear transformation]\label{prob:A5}
Let $\bz$ be a random vector with $\expc[\bz] = \boldsymbol\mu$ and $\var(\bz) = \boldsymbol\Sigma$. For constant $\bA \in \Real^{m\times n}$ and $\bb \in \Real^m$:
\begin{enumerate}[label=(\alph*),nosep]
\item Prove $\var(\bA\bz + \bb) = \bA\boldsymbol\Sigma\bA'$.
\item Use this to derive $\var(\hat{\bb}) = \sigma^2(\bX'\bX)^{-1}$ from $\hat{\bb} - \bb = (\bX'\bX)^{-1}\bX'\be$.
\item Use this to compute $\var(\mathbf{e}) = \sigma^2\bM$ where $\mathbf{e} = \bM\bY$ are the OLS residuals.
\end{enumerate}
\end{exercise}

\begin{solution*}
\textbf{(a)} The constant $\bb$ does not affect variance: $\var(\bA\bz + \bb) = \var(\bA\bz)$. Then directly from the definition,
\begin{align*}
\var(\bA\bz) &= \expc\bigl[(\bA\bz - \bA\boldsymbol\mu)(\bA\bz - \bA\boldsymbol\mu)'\bigr] \\
&= \expc\bigl[\bA(\bz-\boldsymbol\mu)(\bz-\boldsymbol\mu)'\bA'\bigr] = \bA\,\expc\bigl[(\bz-\boldsymbol\mu)(\bz-\boldsymbol\mu)'\bigr]\bA' = \bA\boldsymbol\Sigma\bA'.
\end{align*}

\textbf{(b)} With $\bA = (\bX'\bX)^{-1}\bX'$ (constant given $\bX$) and $\boldsymbol\Sigma = \var(\be) = \sigma^2\bI$:
\[
\var(\hat{\bb}) = \var\bigl((\bX'\bX)^{-1}\bX'\be\bigr) = (\bX'\bX)^{-1}\bX'(\sigma^2\bI)\bX(\bX'\bX)^{-1} = \sigma^2(\bX'\bX)^{-1}.
\]

\textbf{(c)} Note $\mathbf{e} = \bM\bY = \bM(\bX\bb + \be) = \bM\be$ since $\bM\bX = \bzero$. Then
\[
\var(\mathbf{e}) = \bM(\sigma^2\bI)\bM' = \sigma^2\bM\bM = \sigma^2\bM
\]
using symmetry and idempotency of $\bM$. \qed
\end{solution*}

%--- LM-A6 (clm Theorems 3.1, 3.2): Unbiasedness and Variance ---
\begin{exercise}[Unbiasedness and Variance]\label{prob:A6}
(a) Prove $\expc[\hat{\bb}] = \bb$ under (A1)--(A2).\\
(b) Derive $\var(\hat{\bb}) = \sigma^2(\bX'\bX)^{-1}$ under (A1)--(A4).
\end{exercise}

\begin{solution*}
\textbf{(a)} Substituting $\bY = \bX\bb + \be$:
\[
\hat{\bb} = (\bX'\bX)^{-1}\bX'(\bX\bb + \be) = \bb + (\bX'\bX)^{-1}\bX'\be. \tag{$\star$}
\]
Taking expectations and using (A2): $\expc[\hat{\bb}] = \bb + (\bX'\bX)^{-1}\bX'\expc[\be] = \bb$.

\textbf{(b)} By ($\star$), $\hat{\bb}$ is an affine function of $\be$. The variance under a linear transformation $\var(\bA\bz + \bc) = \bA\,\var(\bz)\,\bA'$ (Problem~\ref{prob:A5}) gives, with $\bA = (\bX'\bX)^{-1}\bX'$:
\[
\var(\hat{\bb}) = (\bX'\bX)^{-1}\bX'\,\var(\be)\,\bX(\bX'\bX)^{-1}.
\]
By (A3), $\var(\be) = \sigma^2\bI$. Substituting and using $\bX'\bX(\bX'\bX)^{-1} = \bI$:
\[
\var(\hat{\bb}) = \sigma^2(\bX'\bX)^{-1}\bX'\bX(\bX'\bX)^{-1} = \sigma^2(\bX'\bX)^{-1}. \tag*{\qed}
\]
\end{solution*}

%--- LM-A7 (clm Theorem 3.3): Gauss-Markov ---
\begin{exercise}[Gauss--Markov Theorem]\label{prob:A7}
State and prove the Gauss--Markov theorem: under (A1)--(A4), OLS is BLUE.
\end{exercise}

\begin{solution*}
\emph{Statement.} For any linear unbiased estimator $\tilde{\bb} = \bC\bY$ ($\bC$ a $(k+1)\times n$ constant matrix), $\var(\tilde{\bb}) - \var(\hat{\bb}) \succeq \bzero$.

\textbf{Step 1: Necessary unbiasedness condition.} Under (A1), $\expc[\tilde{\bb}] = \bC\bX\bb = \bb$ for all $\bb \in \Real^{k+1}$ forces $\bC\bX = \bI_{k+1}$ (taking $\bb = e_j$ extracts the $j$-th column).

\textbf{Step 2: Decompose $\bC$.} Write $\bC = (\bX'\bX)^{-1}\bX' + \bD$. Then $\bD\bX = \bC\bX - \bI_{k+1} = \bzero$.

\textbf{Step 3: Compute $\var(\tilde{\bb})$.} Since $\tilde{\bb} = \bC\bY$ and $\var(\bY) = \sigma^2\bI$ by (A1)--(A3):
\begin{align*}
\var(\tilde{\bb}) = \sigma^2\bC\bC' &= \sigma^2[(\bX'\bX)^{-1}\bX' + \bD][\bX(\bX'\bX)^{-1} + \bD'].
\end{align*}
Cross terms vanish: $\bD\bX(\bX'\bX)^{-1} = \bzero$ and $(\bX'\bX)^{-1}\bX'\bD' = (\bD\bX(\bX'\bX)^{-1})' = \bzero$. Hence:
\[
\var(\tilde{\bb}) = \sigma^2(\bX'\bX)^{-1} + \sigma^2\bD\bD' = \var(\hat{\bb}) + \sigma^2\bD\bD'.
\]
$\bD\bD' \succeq \bzero$ (since $\bv'\bD\bD'\bv = \|\bD'\bv\|^2 \geqslant 0$), so $\var(\tilde{\bb}) - \var(\hat{\bb}) \succeq \bzero$. Equality iff $\bD = \bzero$, i.e., $\tilde{\bb} = \hat{\bb}$. \qed
\end{solution*}

%--- LM-A8 (clm Lemma 3.2): Expectation of a quadratic form ---
\begin{exercise}[Expectation of a quadratic form]\label{prob:A8}
Let $\bz$ be a random vector with $\expc[\bz] = \boldsymbol\mu$ and $\var(\bz) = \boldsymbol\Sigma$, and let $\bA$ be a constant symmetric matrix.
\begin{enumerate}[label=(\alph*),nosep]
\item Prove the identity $\expc[\bz'\bA\bz] = \tr(\bA\boldsymbol\Sigma) + \boldsymbol\mu'\bA\boldsymbol\mu$.
\item Apply this with $\bz = \be$, $\bA = \bM$ to derive $\expc[\SSE] = (n-k-1)\sigma^2$, hence $\expc[\MSE] = \sigma^2$.
\end{enumerate}
\end{exercise}

\begin{solution*}
\textbf{(a)} Since $\bz'\bA\bz$ is a scalar, $\bz'\bA\bz = \tr(\bz'\bA\bz) = \tr(\bA\bz\bz')$ by trace cycling. Taking expectations and using linearity of $\tr$:
\[
\expc[\bz'\bA\bz] = \tr\bigl(\bA\,\expc[\bz\bz']\bigr) = \tr\bigl(\bA(\boldsymbol\Sigma + \boldsymbol\mu\boldsymbol\mu')\bigr) = \tr(\bA\boldsymbol\Sigma) + \tr(\bA\boldsymbol\mu\boldsymbol\mu') = \tr(\bA\boldsymbol\Sigma) + \boldsymbol\mu'\bA\boldsymbol\mu,
\]
using $\expc[\bz\bz'] = \boldsymbol\Sigma + \boldsymbol\mu\boldsymbol\mu'$ and trace cycling on the rank-one term.

\textbf{(b)} With $\be$, $\expc[\be] = \bzero$, $\var(\be) = \sigma^2\bI$. Since $\mathbf{e} = \bM\be$, $\SSE = \mathbf{e}'\mathbf{e} = \be'\bM\be$. Apply (a):
\[
\expc[\SSE] = \tr(\bM\cdot\sigma^2\bI) + \bzero = \sigma^2\tr(\bM) = \sigma^2(n-k-1),
\]
using $\tr(\bM) = n-k-1$ from Problem~\ref{prob:A1}. Hence $\expc[\MSE] = \expc[\SSE]/(n-k-1) = \sigma^2$. \qed
\end{solution*}

%--- LM-A9 (clm Theorem 3.4): Unbiasedness of MSE ---
\begin{exercise}[Unbiasedness of MSE]\label{prob:A9}
Prove that $s^2 = \SSE/(n-k-1)$ is an unbiased estimator of $\sigma^2$ under (A1)--(A4).
\end{exercise}

\begin{solution*}
\textbf{Step 1: Express SSE as a quadratic form in $\be$.} Since $\mathbf{e} = \bM\bY = \bM\be$ (as $\bM\bX = \bzero$) and $\bM$ is symmetric idempotent ($\bM'\bM = \bM$):
\[
\SSE = \mathbf{e}'\mathbf{e} = \be'\bM'\bM\be = \be'\bM\be.
\]

\textbf{Step 2: Apply the quadratic-form expectation.} With $\bz = \be$, $\boldsymbol\mu = \bzero$ by (A2), $\boldsymbol\Sigma = \sigma^2\bI$ by (A3), $\bA = \bM$ (Problem~\ref{prob:A8}):
\[
\expc[\SSE] = \tr(\bM \cdot \sigma^2\bI) + \bzero'\bM\bzero = \sigma^2\tr(\bM).
\]

\textbf{Step 3: Evaluate the trace.} By Problem~\ref{prob:A1}, $\tr(\bM) = n-k-1$. Hence:
\[
\expc[\MSE] = \frac{\sigma^2(n-k-1)}{n-k-1} = \sigma^2. \tag*{\qed}
\]
\end{solution*}

%--- LM-A10 (clm Theorem 4.1): MVN linear closure ---
\begin{exercise}[Multivariate normal: closure under linear transformations]\label{prob:A10}
Let $\bz \sim N_n(\boldsymbol\mu, \boldsymbol\Sigma)$. For any constant $\bA \in \Real^{m\times n}$ and $\bb \in \Real^m$, prove
\[
\bA\bz + \bb \sim N_m(\bA\boldsymbol\mu + \bb,\;\bA\boldsymbol\Sigma\bA').
\]
\end{exercise}

\begin{solution*}
Use the moment generating function (MGF) characterization. The MGF of $\bz$ is $M_{\bz}(\bt) = \exp(\bt'\boldsymbol\mu + \tfrac12\bt'\boldsymbol\Sigma\bt)$. For $\bw = \bA\bz + \bb$ and any $\bs \in \Real^m$:
\[
M_{\bw}(\bs) = \expc[\exp(\bs'(\bA\bz+\bb))] = \exp(\bs'\bb)\,\expc[\exp((\bA'\bs)'\bz)] = \exp(\bs'\bb)\,M_{\bz}(\bA'\bs).
\]
Substituting:
\[
M_{\bw}(\bs) = \exp\bigl(\bs'\bb + \bs'\bA\boldsymbol\mu + \tfrac12\bs'\bA\boldsymbol\Sigma\bA'\bs\bigr) = \exp\bigl(\bs'(\bA\boldsymbol\mu+\bb) + \tfrac12\bs'(\bA\boldsymbol\Sigma\bA')\bs\bigr),
\]
which is the MGF of $N_m(\bA\boldsymbol\mu + \bb,\;\bA\boldsymbol\Sigma\bA')$. By uniqueness of MGFs the distribution is identified. \qed
\end{solution*}

%--- LM-A11 (clm Theorem 4.2): Joint normality, uncorrelated <=> independent ---
\begin{exercise}[Joint normality: uncorrelated $\Leftrightarrow$ independent]\label{prob:A11}
Let $(\mathbf{U}', \mathbf{V}')'$ be jointly normal:
\[
\begin{pmatrix} \mathbf{U} \\ \mathbf{V} \end{pmatrix} \sim N\!\left(\begin{pmatrix} \boldsymbol\mu_U \\ \boldsymbol\mu_V \end{pmatrix}, \begin{pmatrix} \boldsymbol\Sigma_{UU} & \boldsymbol\Sigma_{UV} \\ \boldsymbol\Sigma_{VU} & \boldsymbol\Sigma_{VV} \end{pmatrix}\right).
\]
Prove that $\boldsymbol\Sigma_{UV} = \bzero$ if and only if $\mathbf{U}$ and $\mathbf{V}$ are independent.
\end{exercise}

\begin{solution*}
$(\Leftarrow)$: Independence implies uncorrelatedness for any distribution.

$(\Rightarrow)$: With $\boldsymbol\Sigma_{UV} = \bzero$, the covariance is block-diagonal. By Problem~\ref{prob:A10} (extended to the joint MGF),
\[
M_{(\mathbf{U},\mathbf{V})}(\bs, \bt) = \exp\!\left(\bs'\boldsymbol\mu_U + \bt'\boldsymbol\mu_V + \tfrac{1}{2}\bs'\boldsymbol\Sigma_{UU}\bs + \tfrac{1}{2}\bt'\boldsymbol\Sigma_{VV}\bt\right) = M_{\mathbf{U}}(\bs)\,M_{\mathbf{V}}(\bt),
\]
using $\boldsymbol\Sigma_{UV} = \boldsymbol\Sigma_{VU}' = \bzero$ to drop cross terms in the quadratic form. Factorization of the joint MGF as a product of marginal MGFs implies independence. \qed
\end{solution*}

%--- LM-A12 (clm Theorem 4.3): Chi-square of quadratic form [NEW dedicated] ---
\begin{exercise}[Chi-square distribution of quadratic forms]\label{prob:A12}
Let $\bz \sim N_n(\bzero, \bI_n)$ and $\bA$ be an $n\times n$ symmetric idempotent matrix of rank $r$. Prove
\[
\bz'\bA\bz \sim \chi^2_r.
\]
\end{exercise}

\begin{solution*}
\textbf{Step 1: Spectral decomposition.} Since $\bA$ is symmetric, there is an orthogonal $\bP$ with $\bA = \bP\boldsymbol\Lambda\bP'$, where $\boldsymbol\Lambda = \diag(\lambda_1,\ldots,\lambda_n)$. By Problem~\ref{prob:A1}, each $\lambda_i \in \{0,1\}$ and exactly $r = \rank(\bA) = \tr(\bA)$ of them equal $1$. Without loss of generality, the first $r$ are $1$:
\[
\boldsymbol\Lambda = \diag(\underbrace{1,\ldots,1}_{r},\underbrace{0,\ldots,0}_{n-r}).
\]

\textbf{Step 2: Apply orthogonal transformation.} Let $\bw = \bP'\bz$. By Problem~\ref{prob:A10},
\[
\bw \sim N_n(\bP'\bzero,\;\bP'\bI\bP) = N_n(\bzero, \bI_n),
\]
so $w_1,\ldots,w_n$ are i.i.d.\ $N(0,1)$.

\textbf{Step 3: Evaluate the quadratic form.}
\[
\bz'\bA\bz = \bz'\bP\boldsymbol\Lambda\bP'\bz = \bw'\boldsymbol\Lambda\bw = \sum_{i=1}^{n}\lambda_i w_i^2 = \sum_{i=1}^{r}w_i^2.
\]

\textbf{Step 4: Conclude.} Since $w_1,\ldots,w_r$ are i.i.d.\ $N(0,1)$, $\sum_{i=1}^{r}w_i^2 \sim \chi^2_r$. \qed
\end{solution*}

%--- LM-A13 (clm Theorem 4.4): Cochran's theorem ---
\begin{exercise}[Cochran's theorem]\label{prob:A13}
Let $\bz \sim N_n(\bzero, \bI_n)$ and let $\bA_1, \bA_2$ be $n\times n$ symmetric idempotent matrices with $\bA_1\bA_2 = \bzero$.
\begin{enumerate}[label=(\alph*),nosep]
\item Prove that $\bz'\bA_1\bz$ and $\bz'\bA_2\bz$ are independent.
\item Apply (a) to the regression decomposition $\SSE = \be'\bM\be$ vs.\ $\SSR = \be'(\bH-\bH_0)\be$ (under $H_0\colon \beta_1=\cdots=\beta_k=0$, with $\bH_0 = \frac{1}{n}\bone\bone'$) to recover the independence used in the overall $F$ test.
\end{enumerate}
\end{exercise}

\begin{solution*}
\textbf{(a)} Define $\mathbf{U} = \bA_1\bz$ and $\mathbf{V} = \bA_2\bz$. Stack them as a single linear function of $\bz$:
\[
\begin{pmatrix}\mathbf{U}\\\mathbf{V}\end{pmatrix} = \begin{pmatrix}\bA_1 \\ \bA_2\end{pmatrix}\bz.
\]
By Problem~\ref{prob:A10}, $(\mathbf{U}',\mathbf{V}')'$ is \emph{jointly} normal. Their cross-covariance is
\[
\cov(\mathbf{U},\mathbf{V}) = \bA_1\,\var(\bz)\,\bA_2' = \bA_1\bI\bA_2 = \bA_1\bA_2 = \bzero,
\]
using $\bA_2' = \bA_2$ by symmetry. By Problem~\ref{prob:A11}, $\mathbf{U}$ and $\mathbf{V}$ are independent. By symmetric idempotency, $\bz'\bA_1\bz = \|\mathbf{U}\|^2$ depends only on $\mathbf{U}$ and $\bz'\bA_2\bz = \|\mathbf{V}\|^2$ depends only on $\mathbf{V}$, hence the two quadratic forms are independent.

\textbf{(b)} Under $H_0$, $\bY = \beta_0\bone + \be$ and both $(\bH-\bH_0)\bY = (\bH-\bH_0)\be$ and $\bM\bY = \bM\be$ (since $\bH\bone = \bH_0\bone = \bone$ and $\bM\bone = \bzero$). Both $\bH-\bH_0$ and $\bM$ are symmetric idempotent, and $(\bH-\bH_0)\bM = \bH\bM - \bH_0\bM = \bzero - \bzero = \bzero$ (using $\bH\bM = \bzero$ and $\bH_0\bM = \bH_0 - \bH_0\bH = \bzero$). By (a), $\SSR/\sigma^2$ and $\SSE/\sigma^2$ are independent. \qed
\end{solution*}

%--- LM-A14 (clm Theorem 5.2): Distribution of SSE/sigma^2; independence from beta-hat ---
\begin{exercise}[Chi-square distribution and independence of SSE]\label{prob:A14}
Under (A1)--(A5), prove: (a)~$\SSE/\sigma^2 \sim \chi^2_{n-k-1}$. (b)~$\SSE$ is independent of $\hat{\bb}$.
\end{exercise}

\begin{solution*}
\textbf{(a)} Let $\bz = \be/\sigma$. Under (A5), $\bz \sim N(\bzero, \bI_n)$. Then $\SSE/\sigma^2 = \be'\bM\be/\sigma^2 = \bz'\bM\bz$. Since $\bM$ is symmetric idempotent with $\tr(\bM) = n-k-1$, by Problem~\ref{prob:A12}, $\bz'\bM\bz \sim \chi^2_{n-k-1}$.

\textbf{(b)} Stack $\hat{\bb} - \bb = (\bX'\bX)^{-1}\bX'\be$ and $\mathbf{e} = \bM\be$ as a single linear function of $\be$:
\[
\begin{pmatrix}\hat{\bb} - \bb \\ \mathbf{e}\end{pmatrix} = \begin{pmatrix}(\bX'\bX)^{-1}\bX' \\ \bM\end{pmatrix}\be.
\]
Under (A5), $\be$ is normal, so by Problem~\ref{prob:A10} the stacked vector is jointly normal. Their cross-covariance is:
\[
\cov\bigl((\bX'\bX)^{-1}\bX'\be,\;\bM\be\bigr) = (\bX'\bX)^{-1}\bX'(\sigma^2\bI)\bM' = \sigma^2(\bX'\bX)^{-1}\bX'\bM = \bzero,
\]
since $\bX'\bM = (\bM\bX)' = \bzero$. By Problem~\ref{prob:A11}, $\hat{\bb}$ and $\mathbf{e}$ are independent; hence so are $\hat{\bb}$ and $\SSE = \mathbf{e}'\mathbf{e}$. \qed
\end{solution*}

%--- LM-A15 (clm Theorem 6.1, 7.3): ANOVA + overall F test ---
\begin{exercise}[ANOVA and the Overall $F$ Test]\label{prob:A15}
(a) Prove $\SST = \SSR + \SSE$.\\
(b) Under (A1)--(A5) and $H_0\colon \beta_1 = \cdots = \beta_k = 0$, prove $\SSR/\sigma^2 \sim \chi^2_k$ and that $\SSR$ is independent of $\SSE$. Hence derive $F = \MSR/\MSE \sim F_{k, n-k-1}$.
\end{exercise}

\begin{solution*}
\textbf{(a)} Define the centering matrix $\bM_0 = \bI - \frac{1}{n}\bone\bone'$ and $\bH_0 = \frac{1}{n}\bone\bone' = \bI - \bM_0$ (projection onto $\calC(\bone)$). Both are symmetric idempotent, with $\bM_0\bv = \bv - \bar{v}\bone$. Hence $\SST = \bY'\bM_0\bY$, $\SSR = \hat{\bY}'\bM_0\hat{\bY}$ (using $\overline{\hat{Y}} = \bar Y$, which follows from $\bone'\mathbf{e} = 0$), and $\SSE = \mathbf{e}'\mathbf{e}$.

Decompose $\bM_0\bY = \bM_0\hat{\bY} + \bM_0\mathbf{e} = \bM_0\hat{\bY} + \mathbf{e}$, where $\bM_0\mathbf{e} = \mathbf{e}$ because $\bar e = n^{-1}\bone'\mathbf{e} = 0$ (intercept normal equation). Squaring (using $\bM_0^2 = \bM_0$):
\[
\SST = \bY'\bM_0\bY = (\bM_0\bY)'(\bM_0\bY) = \SSR + 2\mathbf{e}'\bM_0\hat{\bY} + \SSE.
\]
The cross term: $\mathbf{e}'\bM_0\hat{\bY} = \mathbf{e}'\hat{\bY} - \bar{Y}\,\mathbf{e}'\bone = 0 - 0 = 0$ (both pieces vanish by the normal equations $\bX'\mathbf{e} = \bzero$). Hence $\SST = \SSR + \SSE$.

\textbf{(b)} Note $\hat{\bY} = \bH\bY$ and $\bar{Y}\bone = \bH_0\bY$, so $\hat{\bY} - \bar{Y}\bone = (\bH - \bH_0)\bY$. Under $H_0$, $\bb = (\beta_0, 0, \ldots, 0)'$, so $\bY = \beta_0\bone + \be$. Since $\bH\bone = \bone$ (as $\bone \in \calC(\bX)$) and $\bH_0\bone = \bone$, $(\bH - \bH_0)\bone = \bzero$, hence $(\bH - \bH_0)\bY = (\bH - \bH_0)\be$. With $\bz = \be/\sigma \sim N(\bzero, \bI)$:
\[
\frac{\SSR}{\sigma^2} = \bz'(\bH - \bH_0)'(\bH - \bH_0)\bz.
\]

\emph{$\bH - \bH_0$ symmetric idempotent of rank $k$:} $\bH\bH_0 = \frac{1}{n}\bH\bone\bone' = \frac{1}{n}\bone\bone' = \bH_0$, and by symmetry $\bH_0\bH = (\bH\bH_0)' = \bH_0$. Then
\[
(\bH - \bH_0)^2 = \bH^2 - \bH\bH_0 - \bH_0\bH + \bH_0^2 = \bH - \bH_0 - \bH_0 + \bH_0 = \bH - \bH_0.
\]
Rank: $\tr(\bH - \bH_0) = (k+1) - 1 = k$. By Problem~\ref{prob:A12}, $\SSR/\sigma^2 = \bz'(\bH - \bH_0)\bz \sim \chi^2_k$.

\emph{Independence:} $(\bH - \bH_0)\bM = \bH\bM - \bH_0\bM = \bzero - (\bH_0 - \bH_0\bH) = \bzero$. By Problem~\ref{prob:A13} (Cochran), $\SSR$ and $\SSE$ are independent.

\emph{$F$ distribution:} $F = \frac{(\SSR/\sigma^2)/k}{(\SSE/\sigma^2)/(n-k-1)} = \MSR/\MSE \sim F_{k,n-k-1}$. \qed
\end{solution*}

\subsection*{Tier B: Apply the Theory}

%--- LM-B1: Leverages ---
\begin{exercise}[Leverages and Residual Variance]\label{prob:B1}
Let $h_{ii}$ denote the $i$-th diagonal entry of $\bH$.\\
(a) Show $0 \leqslant h_{ii} \leqslant 1$. \; (b) Show $\sum_{i=1}^n h_{ii} = k+1$. \; (c) Derive $\var(e_i) = \sigma^2(1-h_{ii})$.
\end{exercise}

\begin{solution*}
\textbf{(a)} Since $\bH$ is symmetric idempotent, $\bH = \bH^2 = \bH'\bH$. Comparing the $(i,i)$ entry, with $h_{ji} = h_{ij}$:
\[
h_{ii} = \sum_j h_{ij}h_{ji} = \sum_j h_{ij}^2 = h_{ii}^2 + \sum_{j \neq i}h_{ij}^2 \geqslant h_{ii}^2,
\]
so $h_{ii}(1-h_{ii}) \geqslant 0$. Combined with $h_{ii} \geqslant 0$ (a sum of squares), this forces $0 \leqslant h_{ii} \leqslant 1$.

\textbf{(b)} $\sum_i h_{ii} = \tr(\bH) = \rank(\bH) = k+1$ (since $\bH$ is symmetric idempotent and projects onto $\calC(\bX)$).

\textbf{(c)} By Problem~\ref{prob:A5}(c), $\var(\mathbf{e}) = \sigma^2\bM$. The $(i,i)$ entry: $\var(e_i) = \sigma^2 m_{ii} = \sigma^2(1 - h_{ii})$. \qed
\end{solution*}

%--- LM-B2: Distribution of SSR under H_0 ---
\begin{exercise}[Distribution of $\SSR$ under $H_0$]\label{prob:B2}
Under (A1)--(A5) and $H_0\colon \beta_1 = \cdots = \beta_k = 0$, prove $\SSR/\sigma^2 \sim \chi^2_k$ and that $\SSR$ is independent of $\SSE$, building only on the lemmas in Tier A.
\end{exercise}

\begin{solution*}
Let $\bH_0 = \frac{1}{n}\bone\bone'$ (projection onto $\calC(\bone)$). Since the model has an intercept, $\bone \in \calC(\bX)$, so $\bH\bone = \bone$ and $\bH\bH_0 = \frac{1}{n}\bH\bone\bone' = \frac{1}{n}\bone\bone' = \bH_0$; by symmetry $\bH_0\bH = \bH_0$. Hence
\[
(\bH - \bH_0)^2 = \bH^2 - \bH\bH_0 - \bH_0\bH + \bH_0^2 = \bH - \bH_0 - \bH_0 + \bH_0 = \bH - \bH_0,
\]
so $\bH-\bH_0$ is symmetric idempotent. Its rank is $\tr(\bH-\bH_0) = (k+1) - 1 = k$ by Problem~\ref{prob:A1} (and the trace identity for $\bH$ from Problem~\ref{prob:A3}).

Under $H_0$, $\bb = (\beta_0, 0, \ldots, 0)'$ so $\bX\bb = \beta_0\bone$ and
\[
\hat{\bY} - \bar Y\bone = (\bH - \bH_0)\bY = (\bH - \bH_0)(\beta_0\bone + \be) = (\bH - \bH_0)\be,
\]
since $(\bH-\bH_0)\bone = \bzero$. With $\bz = \be/\sigma \sim N(\bzero, \bI)$:
\[
\SSR/\sigma^2 = \bz'(\bH-\bH_0)\bz \sim \chi^2_k
\]
by Problem~\ref{prob:A12}.

\emph{Independence:} $(\bH-\bH_0)\bM = \bH\bM - \bH_0\bM = \bzero - (\bH_0 - \bH_0\bH) = \bzero - \bzero = \bzero$. By Problem~\ref{prob:A13} (Cochran), $\SSR$ and $\SSE = \be'\bM\be$ are independent. \qed
\end{solution*}

%--- LM-B3: Testing c'beta = r_0 ---
\begin{exercise}[Testing $\bc'\bb = r_0$]\label{prob:B3}
For a given $\bc \neq \bzero$ and scalar $r_0$, consider $H_0\colon \bc'\bb = r_0$.\\
(a) Derive the distribution of $\bc'\hat{\bb}$. \; (b) Construct the $t$ statistic. \; (c) Show $t^2 \sim F_{1,n-k-1}$.
\end{exercise}

\begin{solution*}
\textbf{(a)} Under (A5), $\hat{\bb} \sim N(\bb, \sigma^2(\bX'\bX)^{-1})$. By Problem~\ref{prob:A10} applied to the row vector $\bc'$: $\bc'\hat{\bb} \sim N(\bc'\bb,\; \sigma^2\bc'(\bX'\bX)^{-1}\bc)$.

\textbf{(b)} Standardize: $Z = \frac{\bc'\hat{\bb} - \bc'\bb}{\sigma\sqrt{\bc'(\bX'\bX)^{-1}\bc}} \sim N(0,1)$. By Problem~\ref{prob:A14}, $V = \SSE/\sigma^2 \sim \chi^2_{n-k-1}$ is independent of $\hat{\bb}$, hence of $Z$. By the definition of $t$, $t = Z/\sqrt{V/(n-k-1)}$:
\[
t = \frac{\bc'\hat{\bb} - r_0}{s\sqrt{\bc'(\bX'\bX)^{-1}\bc}} \sim t_{n-k-1} \quad \text{under } H_0.
\]

\textbf{(c)} $t^2 = \frac{(\bc'\hat{\bb} - r_0)^2}{s^2\,\bc'(\bX'\bX)^{-1}\bc}$. The square of a $t_\nu$ random variable is $F_{1,\nu}$, so $t^2 \sim F_{1,n-k-1}$. \qed
\end{solution*}

%--- LM-B4: Nested F ---
\begin{exercise}[Nested $F$ Test]\label{prob:B4}
Full model with $k$ regressors; reduced model with $k_0 < k$.\\
(a) Prove $\SSE_R \geqslant \SSE_U$. \; (b) Under $H_0$, show $F = \frac{(\SSE_R - \SSE_U)/(k-k_0)}{\SSE_U/(n-k-1)} \sim F_{k-k_0, n-k-1}$.
\end{exercise}

\begin{solution*}
\textbf{(a)} The full model minimizes $\|\bY - \bX\bb\|^2$ over all $\bb \in \Real^{k+1}$. The reduced model minimizes over the subspace $\{\bb : \beta_{k_0+1} = \cdots = \beta_k = 0\}$. The minimum over a subset cannot be smaller, so $\SSE_R \geqslant \SSE_U$.

\textbf{(b)} Let $\bM_U$ and $\bM_R$ be the residual makers for the full and reduced models. By the proof of unbiasedness of MSE, $\SSE_U = \be'\bM_U\be$ and $\SSE_R = \be'\bM_R\be$. Since $\calC(\bX_{\text{red}}) \subseteq \calC(\bX)$, $\bH_U\bH_R = \bH_R\bH_U = \bH_R$, hence:
\begin{itemize}[nosep]
\item $\bM_R - \bM_U$ is symmetric idempotent: $(\bM_R - \bM_U)^2 = \bM_R^2 - \bM_R\bM_U - \bM_U\bM_R + \bM_U^2 = \bM_R - \bM_U$ (using $\bM_R\bM_U = \bM_U = \bM_U\bM_R$, derived from $\bH_U\bH_R = \bH_R$);
\item $\rank(\bM_R - \bM_U) = \tr(\bM_R) - \tr(\bM_U) = (n-k_0-1) - (n-k-1) = k - k_0$;
\item $(\bM_R - \bM_U)\bM_U = \bM_R\bM_U - \bM_U = \bM_U - \bM_U = \bzero$.
\end{itemize}
By Problems~\ref{prob:A12} (chi-square) and~\ref{prob:A13} (Cochran), with $\bz = \be/\sigma$:
\[
\frac{\SSE_R - \SSE_U}{\sigma^2} = \bz'(\bM_R - \bM_U)\bz \sim \chi^2_{k-k_0}, \quad \text{independent of } \frac{\SSE_U}{\sigma^2} \sim \chi^2_{n-k-1}.
\]
Hence $F = \frac{(\SSE_R - \SSE_U)/(k-k_0)}{\SSE_U/(n-k-1)} \sim F_{k-k_0, n-k-1}$. \qed
\end{solution*}

%--- LM-B5: General linear hypothesis F test ---
\begin{exercise}[General linear hypothesis $F$ test]\label{prob:B5}
Consider $H_0\colon \bR\bb = \br$ where $\bR$ is $q\times(k+1)$ with $\rank(\bR) = q$.
\begin{enumerate}[label=(\alph*),nosep]
\item Under $H_0$ and (A5), derive the distribution of $\bR\hat{\bb} - \br$.
\item Show that, under $H_0$,
\[
F = \frac{(\bR\hat{\bb} - \br)'[\bR(\bX'\bX)^{-1}\bR']^{-1}(\bR\hat{\bb} - \br)/q}{\MSE} \sim F_{q,\, n-k-1}.
\]
\item Verify that the individual $t$ test (Problem~\ref{prob:B3}) is the special case $q=1$, $\bR = \bc'$, $\br = r_0$, with $F = t^2$.
\end{enumerate}
\end{exercise}

\begin{solution*}
\textbf{(a)} By Problem~\ref{prob:A10}, $\hat{\bb} = \bb + (\bX'\bX)^{-1}\bX'\be \sim N\bigl(\bb,\;\sigma^2(\bX'\bX)^{-1}\bigr)$ under (A5). Applying the same lemma again to $\bR\hat\bb$:
\[
\bR\hat{\bb} \sim N\bigl(\bR\bb,\;\sigma^2\bR(\bX'\bX)^{-1}\bR'\bigr).
\]
Under $H_0$: $\bR\hat{\bb} - \br \sim N(\bzero,\;\sigma^2\bR(\bX'\bX)^{-1}\bR')$. The covariance is positive definite because $\rank(\bR) = q$ and $(\bX'\bX)^{-1} \succ 0$.

\textbf{(b)} Let $\boldsymbol\Sigma = \bR(\bX'\bX)^{-1}\bR'$ (positive definite, $q\times q$). Standardize: $\boldsymbol\Sigma^{-1/2}(\bR\hat{\bb} - \br)/\sigma \sim N(\bzero, \bI_q)$, so by Problem~\ref{prob:A12},
\[
\frac{(\bR\hat{\bb} - \br)'\boldsymbol\Sigma^{-1}(\bR\hat{\bb} - \br)}{\sigma^2} \sim \chi^2_q.
\]
By Problem~\ref{prob:A14}, $\SSE/\sigma^2 \sim \chi^2_{n-k-1}$ and is independent of $\hat{\bb}$, hence independent of $\bR\hat{\bb}$. Forming the ratio:
\[
F = \frac{[(\bR\hat{\bb} - \br)'\boldsymbol\Sigma^{-1}(\bR\hat{\bb} - \br)/\sigma^2]/q}{[\SSE/\sigma^2]/(n-k-1)} = \frac{(\bR\hat{\bb} - \br)'\boldsymbol\Sigma^{-1}(\bR\hat{\bb} - \br)/q}{\MSE} \sim F_{q,n-k-1}.
\]

\textbf{(c)} With $q=1$, $\bR = \bc'$, $\br = r_0$: $\boldsymbol\Sigma = \bc'(\bX'\bX)^{-1}\bc$ (a scalar) and the numerator is $(\bc'\hat{\bb} - r_0)^2/[\bc'(\bX'\bX)^{-1}\bc]$. Then
\[
F = \frac{(\bc'\hat{\bb} - r_0)^2}{\MSE\cdot\bc'(\bX'\bX)^{-1}\bc} = \left(\frac{\bc'\hat{\bb} - r_0}{\SE(\bc'\hat{\bb})}\right)^2 = t^2,
\]
matching Problem~\ref{prob:B3} ($F_{1,n-k-1}$ equals the square of $t_{n-k-1}$). \qed
\end{solution*}

%--- LM-B6: GLS ---
\begin{exercise}[GLS is BLUE]\label{prob:B6}
Suppose $\var(\be) = \sigma^2\boldsymbol\Omega$ where $\boldsymbol\Omega$ is known, positive definite, $\neq \bI$.\\
(a) Show OLS is still unbiased. Derive its actual variance.\\
(b) Define $\hat{\bb}_{\mathrm{GLS}} = (\bX'\boldsymbol\Omega^{-1}\bX)^{-1}\bX'\boldsymbol\Omega^{-1}\bY$. Prove it is BLUE.
\end{exercise}

\begin{solution*}
\textbf{(a)} \emph{Unbiasedness:} $\expc[\hat{\bb}] = \bb + (\bX'\bX)^{-1}\bX'\expc[\be] = \bb$ (uses only (A1)--(A2), not (A3)).

\emph{Actual variance:}
\[
\var(\hat{\bb}) = (\bX'\bX)^{-1}\bX'\,\var(\be)\,\bX(\bX'\bX)^{-1} = \sigma^2(\bX'\bX)^{-1}\bX'\boldsymbol\Omega\bX(\bX'\bX)^{-1}.
\]
This differs from $\sigma^2(\bX'\bX)^{-1}$ when $\boldsymbol\Omega \neq \bI$, so the usual OLS standard errors are incorrect.

\textbf{(b)} Factor $\boldsymbol\Omega = \bP\bP'$ (Cholesky or symmetric square root). Pre-multiplying the model by $\bP^{-1}$:
\[
\bP^{-1}\bY = \bP^{-1}\bX\bb + \bP^{-1}\be.
\]
The transformed errors satisfy $\var(\bP^{-1}\be) = \bP^{-1}(\sigma^2\boldsymbol\Omega)(\bP^{-1})' = \sigma^2\bI$. The transformed model thus satisfies (A1)--(A4) with spherical errors. OLS on the transformed model gives:
\[
\hat{\bb}_{\mathrm{GLS}} = [(\bP^{-1}\bX)'(\bP^{-1}\bX)]^{-1}(\bP^{-1}\bX)'(\bP^{-1}\bY) = (\bX'\boldsymbol\Omega^{-1}\bX)^{-1}\bX'\boldsymbol\Omega^{-1}\bY.
\]
By Gauss--Markov on the transformed model, $\hat{\bb}_{\mathrm{GLS}}$ is BLUE. \qed
\end{solution*}

%=============================================================================
%\newpage
%\begin{center}
%\rule{0.6\textwidth}{0.4pt}\\[2mm]
%{\Large\bfseries Appendix: Key Derivations Checklist}\\[1mm]
%\rule{0.6\textwidth}{0.4pt}
%\end{center}
%
%\noindent\textbf{Matrix Calculus}
%\begin{enumerate}[label=\textbf{MC\arabic*.},leftmargin=2cm,nosep]
%\item $\tr A'B = \tr BA'$ (trace cycling).
%\item $x'Ax = 0$ for all $x$ $\Leftrightarrow$ $A$ skew-symmetric; $+$ symmetry $\Rightarrow$ $A = 0$.
%\item $(A \otimes B)(C \otimes D) = AC \otimes BD$.
%\item $\avec(ABC) = (C' \otimes A)\avec B$.
%\item $\dd|X| = |X|\tr X^{-1}\dd X$; \; $\dd\log|X| = \tr X^{-1}\dd X$.
%\item $\dd X^{-1} = -X^{-1}(\dd X)X^{-1}$.
%\item Derivatives: $a'x$, $x'Ax$, $\tr X'AX$, $\log|X'X|$, $\tr X^k$, $AX^{-1}B$.
%\item Hessian of $x'Ax$: $Hf = A+A'$; symmetrization rule $(B+B')/2$.
%\item Quadratic optimization: unconstrained ($\hat{x} = A^{-1}b$) and constrained ($c'x = 0$).
%\item MLE for multivariate normal: $\hat\mu = \bar{y}$, $\hat\Omega = \frac{1}{n}\sum(y_i-\bar{y})(y_i-\bar{y})'$.
%\end{enumerate}
%
%\medskip
%\noindent\textbf{Multiple Linear Regression}
%\begin{enumerate}[label=\textbf{LM\arabic*.},leftmargin=2cm,nosep]
%\item OLS estimator via differentiation and projection.
%\item $\bH$, $\bM$: symmetric idempotent, $\bH\bX=\bX$, $\bM\bX=\bzero$, traces.
%\item Unbiasedness of $\hat\bb$ and $\var(\hat\bb) = \sigma^2(\bX'\bX)^{-1}$.
%\item Gauss--Markov theorem (BLUE).
%\item Unbiasedness of $\MSE$ via trace identity.
%\item $\SSE/\sigma^2 \sim \chi^2_{n-k-1}$ and independence from $\hat\bb$.
%\item $\SST = \SSR + \SSE$.
%\item Overall $F$ test via $\bH - \bH_0$ idempotent.
%\item Individual $t$ test.
%\item General $F$ test for $\bR\bb = \br$.
%\end{enumerate}

\end{document}
